\chapter{约束最优化理论}\label{chap:constrained-theory}
% 第6章第1节：引言与基本理论
\section{约束最优化问题概述}\label{sec:ch6-intro}

\subsection{问题描述}

一般约束最优化问题（非线性规划，NLP）的标准形式：
\begin{align}
\min \quad & f(x) \notag \\
\text{s.t.} \quad & g_i(x) \leq 0, \quad i = 1, 2, \ldots, m \label{eq:nlp} \\
& h_j(x) = 0, \quad j = 1, 2, \ldots, l \notag
\end{align}

其中 $x \in \mathbb{R}^n$，$f: \mathbb{R}^n \to \mathbb{R}$ 为目标函数，$g_i$ 为不等式约束，$h_j$ 为等式约束。

\begin{definition}[可行集]
可行集（Feasible Set）定义为：
\[
\mathcal{F} = \{x \in \mathbb{R}^n : g_i(x) \leq 0, \, \forall i; \, h_j(x) = 0, \, \forall j\}
\]
\end{definition}

\begin{definition}[积极约束与有效集]
在点 $\bar{x} \in \mathcal{F}$ 处：
\begin{itemize}
    \item 若 $g_i(\bar{x}) = 0$，称约束 $g_i$ 在 $\bar{x}$ 处\textbf{积极的}（active）
    \item 有效集（Active Set）：$\mathcal{A}(\bar{x}) = \{i : g_i(\bar{x}) = 0\} \cup \{m+j : j=1,\ldots,l\}$
\end{itemize}
\end{definition}

\subsection{可行方向}\label{sec:ch6-feasible-direction}

\begin{definition}[可行方向]
设 $\bar{x} \in \mathcal{F}$，向量 $d \in \mathbb{R}^n$ 称为 $\bar{x}$ 处的\textbf{可行方向}，如果存在 $\bar{\alpha} > 0$，使得：
\[
\bar{x} + \alpha d \in \mathcal{F}, \quad \forall \alpha \in [0, \bar{\alpha}]
\]
\end{definition}

\begin{conceptbox}[下降可行方向]
对于问题 \eqref{eq:nlp}，在点 $\bar{x}$ 处的可行方向 $d$ 还需满足下降条件。若 $d$ 满足：
\begin{align*}
\nabla f(\bar{x})^T d &< 0 \quad \text{(下降方向)} \\
\nabla g_i(\bar{x})^T d &\leq 0, \quad i \in \mathcal{A}(\bar{x}) \quad \text{(不等式约束)} \\
\nabla h_j(\bar{x})^T d &= 0, \quad j = 1, \ldots, l \quad \text{(等式约束)}
\end{align*}
则 $d$ 称为\textbf{下降可行方向}。
\end{conceptbox}

\begin{example}[可行方向验证]
考虑问题：
\[
\min \, x_1^2 + x_2^2, \quad \text{s.t. } x_1 + x_2 \leq 2, \, x_1 \geq 0, \, x_2 \geq 0
\]

\textbf{解：}

在点 $\bar{x} = (1, 1)$ 处，$g_1(x) = x_1 + x_2 - 2$ 是积极约束。

$\nabla g_1 = (1, 1)$，$\nabla f = (2x_1, 2x_2) = (2, 2)$。

取 $d = (-1, -1)$：
\begin{align*}
\nabla f^T d &= (2, 2) \cdot (-1, -1) = -4 < 0 \quad \text{(下降)} \\
\nabla g_1^T d &= (1, 1) \cdot (-1, -1) = -2 \leq 0 \quad \text{(可行)}
\end{align*}
所以 $d = (-1, -1)$ 是下降可行方向。
\end{example}

\subsection{切锥与线性化可行方向}\label{sec:tangent-cone}

\begin{definition}[切锥（Tangent Cone）]
设 $\mathcal{F} \subseteq \mathbb{R}^n$，$\bar{x} \in \mathcal{F}$。切锥 $T_{\mathcal{F}}(\bar{x})$ 定义为所有满足以下条件的向量 $d$ 的集合：存在序列 $\{x_k\} \subset \mathcal{F}$，$x_k \to \bar{x}$，以及 $\{\alpha_k\} \subset \mathbb{R}_+$，$\alpha_k \to 0$，使得：
\[
d = \lim_{k \to \infty} \frac{x_k - \bar{x}}{\alpha_k}
\]
\end{definition}

\begin{definition}[线性化可行方向（LFD）]
在点 $\bar{x}$ 处的线性化可行方向集为：
\[
\mathcal{L}(\bar{x}) = \{d \in \mathbb{R}^n : \nabla g_i(\bar{x})^T d \leq 0, \, i \in \mathcal{A}(\bar{x}); \, \nabla h_j(\bar{x})^T d = 0, \, j = 1,\ldots,l\}
\]
\end{definition}

\begin{notebox}[切锥与LFD的关系]
一般来说，$T_{\mathcal{F}}(\bar{x}) \subseteq \mathcal{L}(\bar{x})$。当约束满足某些\textbf{约束规范}（Constraint Qualification）时，二者相等。约束规范保证了KKT条件的必要性。
\end{notebox}

\subsection{Farkas引理}\label{sec:farkas}

Farkas引理是约束优化理论的基石，用于证明KKT条件的必要性。

\begin{theorem}[Farkas引理]
设 $A \in \mathbb{R}^{m \times n}$，$b \in \mathbb{R}^m$。则以下两个系统有且仅有一个有解：

\textbf{系统1}：$Ax \leq 0$，$b^T x > 0$

\textbf{系统2}：$A^T y = b$，$y \geq 0$
\end{theorem}

等价表述：
\begin{kktbox}[Farkas引理（标准形式）]
设 $a_1, \ldots, a_m, b \in \mathbb{R}^n$。则以下两个系统有且仅有一个有解：

\textbf{系统1}：$b^T d > 0$ 且 $a_i^T d \leq 0$，$i = 1, \ldots, m$

\textbf{系统2}：$b = \sum_{i=1}^m \lambda_i a_i$，$\lambda_i \geq 0$
\end{kktbox}

\textbf{几何解释}：要么 $b$ 在锥 $\text{cone}\{a_1, \ldots, a_m\}$ 内（系统2），要么存在超平面将 $b$ 与该锥分离（系统1）。

\begin{corollary}[Gordan定理]
以下两个系统有且仅有一个有解：

\textbf{系统1}：$Ax < 0$

\textbf{系统2}：$A^T y = 0$，$y \geq 0$，$y \neq 0$
\end{corollary}

\subsection{无约束问题最优性条件回顾}

在深入约束问题之前，先回顾无约束问题的最优性条件。

\begin{theorem}[无约束一阶必要条件]
设 $f$ 连续可微。若 $x^*$ 是局部极小点，则 $\nabla f(x^*) = 0$。
\end{theorem}

\begin{theorem}[无约束二阶必要条件]
设 $f$ 二阶连续可微。若 $x^*$ 是局部极小点，则 $\nabla^2 f(x^*) \succeq 0$（半正定）。
\end{theorem}

\begin{theorem}[无约束二阶充分条件]
设 $f$ 二阶连续可微。若 $\nabla f(x^*) = 0$ 且 $\nabla^2 f(x^*) \succ 0$（正定），则 $x^*$ 是严格局部极小点。
\end{theorem}

% 第6章第2节：KKT一阶最优性条件
\section{一阶最优性条件（KKT条件）}\label{sec:ch6-kkt}

\subsection{等式约束问题}

先考虑只有等式约束的问题：
\[
\min \, f(x), \quad \text{s.t. } h_j(x) = 0, \quad j = 1, \ldots, l
\]

\begin{theorem}[等式约束一阶必要条件]
设 $f$ 和 $h_j$ 连续可微，$x^*$ 是局部极小点，且 $\{\nabla h_j(x^*)\}_{j=1}^l$ 线性无关（LICQ）。则存在 $\mu^* \in \mathbb{R}^l$，使得：
\begin{align}
\nabla f(x^*) + \sum_{j=1}^l \mu_j^* \nabla h_j(x^*) &= 0 \label{eq:eq-lagrange} \\
h_j(x^*) &= 0, \quad j = 1, \ldots, l
\end{align}
\end{theorem}

\begin{conceptbox}[拉格朗日函数]
定义拉格朗日函数：
\[
\mathcal{L}(x, \mu) = f(x) + \sum_{j=1}^l \mu_j h_j(x)
\]
则条件 \eqref{eq:eq-lagrange} 等价于 $\nabla_x \mathcal{L}(x^*, \mu^*) = 0$。
\end{conceptbox}

\begin{example}[等式约束问题]
\[
\min \, x_1^2 + x_2^2, \quad \text{s.t. } x_1 + x_2 = 1
\]

\textbf{解：}

$\mathcal{L} = x_1^2 + x_2^2 + \mu(x_1 + x_2 - 1)$

KKT条件：
\begin{align*}
\frac{\partial \mathcal{L}}{\partial x_1} &= 2x_1 + \mu = 0 \\
\frac{\partial \mathcal{L}}{\partial x_2} &= 2x_2 + \mu = 0 \\
x_1 + x_2 &= 1
\end{align*}

解得：$x_1 = x_2 = 0.5$，$\mu = -1$。最优值 $f^* = 0.5$。
\end{example}

\subsection{一般约束问题的KKT条件}\label{sec:kkt-general}

\begin{kktbox}[KKT一阶必要条件]
考虑问题 \eqref{eq:nlp}。设 $f, g_i, h_j$ 连续可微，$x^*$ 是局部极小点，且某约束规范（CQ）成立。则存在 $\lambda^* \in \mathbb{R}^m$，$\mu^* \in \mathbb{R}^l$，满足：

\textbf{(1) 平稳性（Stationarity）}：
\[
\nabla f(x^*) + \sum_{i=1}^m \lambda_i^* \nabla g_i(x^*) + \sum_{j=1}^l \mu_j^* \nabla h_j(x^*) = 0
\]

\textbf{(2) 原始可行性（Primal Feasibility）}：
\[
g_i(x^*) \leq 0 \, (i=1,\ldots,m), \quad h_j(x^*) = 0 \, (j=1,\ldots,l)
\]

\textbf{(3) 对偶可行性（Dual Feasibility）}：
\[
\lambda_i^* \geq 0, \quad i = 1, \ldots, m
\]

\textbf{(4) 互补松弛条件（Complementary Slackness）}：
\[
\lambda_i^* \cdot g_i(x^*) = 0, \quad i = 1, \ldots, m
\]
\end{kktbox}

\begin{conceptbox}[KKT条件的直观理解]
\begin{itemize}
    \item \textbf{平稳性}：在最优解处，目标函数的梯度可由约束梯度的线性组合``抵消''
    \item \textbf{互补松弛}：若 $\lambda_i^* > 0$，则 $g_i(x^*) = 0$（约束积极）；若 $g_i(x^*) < 0$，则 $\lambda_i^* = 0$
    \item $\lambda_i^*$ 的大小反映约束的``重要性''——影子价格
\end{itemize}
\end{conceptbox}

\begin{figure}[H]
\centering
\begin{tikzpicture}[scale=0.85]
    % 坐标轴
    \draw[->,thick] (-0.5,0) -- (4.5,0) node[right] {$x_1$};
    \draw[->,thick] (0,-0.5) -- (0,3.8) node[above] {$x_2$};
    % 可行域（三角形）
    \fill[green!12] (0,0) -- (3,0) -- (0,3) -- cycle;
    \draw[thick,blue] (0,0) -- (3,0) -- (0,3) -- cycle;
    % 约束边界标注
    \node[blue,scale=0.8] at (2.5,0.25) {\small $g_2: x_1+x_2=3$};
    \node[blue,scale=0.8] at (-0.35,1.8) {\small $g_3: x_1=0$};
    \node[blue,scale=0.8] at (1.6,-0.35) {\small $g_4: x_2=0$};
    % 目标函数等值线（以(3,2)为圆心的同心圆）
    \foreach \r in {0.8,1.6,2.4} {
        \draw[gray!60,dashed] (3,2) circle (\r);
    }
    \node[gray,scale=0.7] at (3.7,3.2) {\small $f$等值线};
    % 最优点（在约束 g_2 = x_1+x_2 = 3 上）
    \fill[red] (1.5,1.5) circle (3.5pt) node[above right,scale=0.85] {\small $x^*\!=\!(1.5,\!1.5)$};
    % 目标函数梯度 ∇f = 2(x-3, y-2) = 2(-1.5, -0.5) = (-3, -1)
    \draw[->,very thick,red!70!black] (1.5,1.5) -- (0.8,1.27) node[left,scale=0.8] {$\nabla f(x^*)$};
    % 约束 g_2 的梯度 ∇g_2 = (1,1)
    \draw[->,very thick,blue!70!black] (1.5,1.5) -- (2.2,2.2) node[above,scale=0.8] {$\nabla g_2(x^*)$};
    % KKT平衡：-∇f 应与 ∇g_2 同方向
    \draw[->,very thick,orange!80!black] (1.5,1.5) -- (2.2,1.73) node[right,scale=0.8] {$-\nabla f(x^*)$};
    % 说明文字
    \node[scale=0.75] at (2.2,3.5) {\small 在最优解处：$-\nabla f = \lambda_2 \nabla g_2$};
\end{tikzpicture}
\caption{KKT条件的几何解释：在最优解 $x^*$ 处，负目标函数梯度 $-\nabla f$ 可由积极约束梯度的非负线性组合表示。本例中 $-\nabla f = 3(1,1) = 3\nabla g_2$，故 $\lambda_2 = 3 > 0$}
\label{fig:kkt-intuitive}
\end{figure}

\begin{example}[KKT条件求解]\label{ex:kkt-solve}
求解下列约束优化问题
\[
\min \, (x_1 - 2)^2 + (x_2 - 1)^2, \quad \text{s.t. } x_1 + x_2 \leq 2, \; x_1, x_2 \geq 0
\]

\textbf{解：}

\noindent 将约束写为标准形式：$g_1(x) = x_1 + x_2 - 2 \leq 0$，$g_2(x) = -x_1 \leq 0$，$g_3(x) = -x_2 \leq 0$。

梯度：$\nabla f(x) = \begin{bmatrix} 2(x_1-2) \\ 2(x_2-1) \end{bmatrix}$，$\nabla g_1 = \begin{bmatrix} 1 \\ 1 \end{bmatrix}$，$\nabla g_2 = \begin{bmatrix} -1 \\ 0 \end{bmatrix}$，$\nabla g_3 = \begin{bmatrix} 0 \\ -1 \end{bmatrix}$。

列写 KKT 条件的四个性质：
\begin{description}
\item[\textbf{\textcircled{1} 平稳性}（Stationarity）]
$2(x_1-2) + \lambda_1 - \lambda_2 = 0$，\quad $2(x_2-1) + \lambda_1 - \lambda_3 = 0$
\item[\textbf{\textcircled{2} 原始可行性}（Primal Feasibility）]
$x_1 + x_2 \leq 2$，$x_1 \geq 0$，$x_2 \geq 0$
\item[\textbf{\textcircled{3} 对偶可行性}（Dual Feasibility）]
$\lambda_1 \geq 0$，$\lambda_2 \geq 0$，$\lambda_3 \geq 0$
\item[\textbf{\textcircled{4} 互补松弛性}（Complementary Slackness）]
$\lambda_1(x_1 + x_2 - 2) = 0$，\quad
$\lambda_2(-x_1) = 0$，\quad
$\lambda_3(-x_2) = 0$
\end{description}

由互补松弛性（\textcircled{4}），按积极约束的组合分情形讨论（共 $2^3 = 8$ 种）：

\medskip\noindent\textbf{Case 1：}$g_1$ 积极，$g_2,g_3$ 不积极

$\lambda_2 \stackrel{\textcircled{4}}{=} 0$，$\lambda_3 \stackrel{\textcircled{4}}{=} 0$，$x_1 + x_2 \stackrel{\textcircled{4}}{=} 2$。

平稳性（\textcircled{1}） 化为：
\[
2(x_1-2) + \lambda_1 = 0,\quad 2(x_2-1) + \lambda_1 = 0
\]
两式相减得 $2(x_1 - x_2 - 1) = 0 \Rightarrow x_1 - x_2 = 1$。

联立 $x_1 + x_2 = 2$，解得 $x_1 = \frac{3}{2}$，$x_2 = \frac{1}{2}$，$\lambda_1 = 1$。

检验：$\lambda_1 = 1 > 0$ 满足对偶可行性 （\textcircled{3}）；$x_1,x_2 > 0$ 满足原始可行性 （\textcircled{2}）。

$\Rightarrow$ \textbf{KKT点}：$x^* = \left(\frac{3}{2},\, \frac{1}{2}\right)$，$\lambda^* = (1, 0, 0)$。

\bigskip\noindent\textbf{Case 2：}$g_2$ 积极，$g_1,g_3$ 不积极

$\lambda_1 \stackrel{\textcircled{4}}{=} 0$，$\lambda_3 \stackrel{\textcircled{4}}{=} 0$，$x_1 \stackrel{\textcircled{4}}{=} 0$。

平稳性（\textcircled{1}）：$2(0-2) - \lambda_2 = 0 \Rightarrow \lambda_2 = -4$。

$\lambda_2 = -4 < 0$，违反对偶可行性 （\textcircled{3}）。$\Rightarrow$ \textbf{不是 KKT 点}。

\bigskip\noindent\textbf{Case 3：}$g_3$ 积极，$g_1,g_2$ 不积极

$\lambda_1 \stackrel{\textcircled{4}}{=} 0$，$\lambda_2 \stackrel{\textcircled{4}}{=} 0$，$x_2 \stackrel{\textcircled{4}}{=} 0$。

平稳性（\textcircled{1}）：$2(x_1-2) = 0 \Rightarrow x_1 = 2$；$2(0-1) - \lambda_3 = 0 \Rightarrow \lambda_3 = -2$。

$\lambda_3 = -2 < 0$，违反对偶可行性 （\textcircled{3}）。$\Rightarrow$ \textbf{不是 KKT 点}。

\bigskip\noindent\textbf{Case 4：}$g_1,g_2$ 积极，$g_3$ 不积极

$\lambda_3 \stackrel{\textcircled{4}}{=} 0$，$x_1 + x_2 = 2$，$x_1 = 0 \Rightarrow x_2 = 2$。

平稳性（\textcircled{1}）：$2(0-2) + \lambda_1 - \lambda_2 = 0 \Rightarrow \lambda_1 - \lambda_2 = 4$；$2(2-1) + \lambda_1 = 0 \Rightarrow \lambda_1 = -2$。

$\lambda_1 = -2 < 0$，违反对偶可行性 （\textcircled{3}）。$\Rightarrow$ \textbf{不是 KKT 点}。

\bigskip\noindent\textbf{Case 5：}$g_1,g_3$ 积极，$g_2$ 不积极

$\lambda_2 \stackrel{\textcircled{4}}{=} 0$，$x_1 + x_2 = 2$，$x_2 = 0 \Rightarrow x_1 = 2$。

平稳性（\textcircled{1}）：$2(2-2) + \lambda_1 = 0 \Rightarrow \lambda_1 = 0$；$2(0-1) + 0 - \lambda_3 = 0 \Rightarrow \lambda_3 = -2$。

$\lambda_3 = -2 < 0$，违反对偶可行性 （\textcircled{3}）。$\Rightarrow$ \textbf{不是 KKT 点}。

\bigskip\noindent\textbf{Case 6：}$g_2,g_3$ 积极，$g_1$ 不积极

$\lambda_1 \stackrel{\textcircled{4}}{=} 0$，$x_1 = 0$，$x_2 = 0$。

平稳性（\textcircled{1}）：$2(0-2) - \lambda_2 = 0 \Rightarrow \lambda_2 = -4$。

$\lambda_2 = -4 < 0$，违反对偶可行性 （\textcircled{3}）。$\Rightarrow$ \textbf{不是 KKT 点}。

\bigskip\noindent\textbf{Case 7：}$g_1,g_2,g_3$ 均积极

$x_1 + x_2 = 2$，$x_1 = 0$，$x_2 = 0$ $\Rightarrow$ $0 = 2$，矛盾。

$\Rightarrow$ \textbf{不可行}，无解。

\bigskip\noindent\textbf{Case 8：}$g_1,g_2,g_3$ 均不积极

$\lambda_1 \stackrel{\textcircled{4}}{=} 0$，$\lambda_2 \stackrel{\textcircled{4}}{=} 0$，$\lambda_3 \stackrel{\textcircled{4}}{=} 0$。

平稳性（\textcircled{1}）：$2(x_1-2) = 0 \Rightarrow x_1 = 2$；$2(x_2-1) = 0 \Rightarrow x_2 = 1$。

但 $x_1 + x_2 = 3 > 2$，违反原始可行性 （\textcircled{2}）。$\Rightarrow$ \textbf{不是可行点}。

\medskip\noindent\rule{\textwidth}{0.4pt}

综上，仅 Case 1 满足全部 KKT 条件。问题为凸优化（目标函数凸，约束凸），由 KKT 充分性定理，$x^* = \left(\frac{3}{2},\,\frac{1}{2}\right)$ 是全局最优解，最优值 $f^* = \frac{1}{2}$。
\end{example}

\subsection{约束规范（Constraint Qualifications）}\label{sec:cq}

KKT条件的成立需要约束规范。常用的约束规范：

\begin{definition}[LICQ — 线性无关约束规范]
在 $x^*$ 处，积极约束的梯度 $\{\nabla g_i(x^*) : i \in \mathcal{A}(x^*)\} \cup \{\nabla h_j(x^*)\}$ 线性无关。
\end{definition}

\begin{definition}[MFCQ — Mangasarian-Fromovitz约束规范]
存在向量 $d \in \mathbb{R}^n$，使得：
\begin{align*}
\nabla g_i(x^*)^T d &< 0, \quad i \in \mathcal{A}(x^*) \\
\nabla h_j(x^*)^T d &= 0, \quad j = 1, \ldots, l
\end{align*}
且 $\{\nabla h_j(x^*)\}$ 线性无关。
\end{definition}

\begin{definition}[Slater条件]
对于凸优化问题，若存在严格可行点 $\tilde{x}$ 使得：
\[
g_i(\tilde{x}) < 0 \; (\forall i), \quad h_j(\tilde{x}) = 0 \; (\forall j)
\]
且等式约束为仿射，则Slater条件成立。
\end{definition}

\begin{theorem}[CQ的蕴含关系]
对于凸优化问题：
\[
\text{LICQ} \Rightarrow \text{MFCQ}; \quad \text{Slater (凸)} \Rightarrow \text{KKT 必要}
\]
Slater条件是最常用的凸优化CQ。
\end{theorem}

\begin{example}[Slater条件验证]
\[
\min \, x_1^2 + x_2^2, \quad \text{s.t. } x_1 + x_2 \leq 2, \; x_1 \geq 0, \; x_2 \geq 0
\]

\textbf{解：}

取 $\tilde{x} = (0.5, 0.5)$：
\begin{align*}
g_1(0.5, 0.5) &= 1 - 2 = -1 < 0 \\
g_2(0.5, 0.5) &= -0.5 < 0 \\
g_3(0.5, 0.5) &= -0.5 < 0
\end{align*}
Slater条件满足，KKT条件是最优性的充要条件。
\end{example}

\subsection{凸优化的KKT充分性}

\begin{theorem}[KKT充分条件]\label{thm:kkt-sufficient}
若问题 \eqref{eq:nlp} 是凸优化问题（$f$ 凸，$g_i$ 凸，$h_j$ 仿射），且 $(x^*, \lambda^*, \mu^*)$ 满足KKT条件，则 $x^*$ 是全局最优解。
\end{theorem}

\textbf{证明概要}：
\begin{align*}
f(x) &\geq f(x^*) + \nabla f(x^*)^T(x - x^*) \\
&= f(x^*) - \sum_i \lambda_i^* \nabla g_i(x^*)^T(x - x^*) - \sum_j \mu_j^* \nabla h_j(x^*)^T(x - x^*)
\end{align*}
利用凸性 $g_i(x) \geq g_i(x^*) + \nabla g_i(x^*)^T(x-x^*)$ 和互补松弛条件即可得证。

\subsection{拉格朗日函数与鞍点}

\begin{definition}[一般拉格朗日函数]
\[
\mathcal{L}(x, \lambda, \mu) = f(x) + \sum_{i=1}^m \lambda_i g_i(x) + \sum_{j=1}^l \mu_j h_j(x)
\]
其中 $\lambda \geq 0$。
\end{definition}

\begin{definition}[鞍点]
$(x^*, \lambda^*, \mu^*)$ 称为 $\mathcal{L}$ 的鞍点，如果：
\[
\mathcal{L}(x^*, \lambda, \mu) \leq \mathcal{L}(x^*, \lambda^*, \mu^*) \leq \mathcal{L}(x, \lambda^*, \mu^*), \quad \forall x, \forall \lambda \geq 0, \forall \mu
\]
\end{definition}

\begin{theorem}[鞍点定理]
若 $(x^*, \lambda^*, \mu^*)$ 是 $\mathcal{L}$ 的鞍点，则 $x^*$ 是原问题最优解，$(\lambda^*, \mu^*)$ 是对偶问题最优解，且无对偶间隙。
\end{theorem}

% 第6章第3节：二阶最优性条件
\section{二阶最优性条件}\label{sec:second-order}

\subsection{临界锥（Critical Cone）}

二阶条件需要限制在临界锥上，因为积极约束之外的方向不影响局部最优性。

\begin{definition}[临界锥]
设 $x^*$ 是满足KKT条件的点，对应乘子 $(\lambda^*, \mu^*)$。临界锥定义为：
\[
\mathcal{C}(x^*) = \left\{d \in \mathbb{R}^n :
\begin{array}{l}
\nabla g_i(x^*)^T d = 0, \; i \in \mathcal{A}(x^*) \text{ 且 } \lambda_i^* > 0 \\
\nabla g_i(x^*)^T d \leq 0, \; i \in \mathcal{A}(x^*) \text{ 且 } \lambda_i^* = 0 \\
\nabla h_j(x^*)^T d = 0, \; j = 1, \ldots, l
\end{array}\right\}
\]
\end{definition}

\begin{notebox}[临界锥的含义]
\begin{itemize}
    \item 对于强积极约束（$\lambda_i^* > 0$）：方向必须与之正交（不能进一步违反）
    \item 对于弱积极约束（$\lambda_i^* = 0$）：方向可以非正交（约束仍然允许移动）
    \item 临界锥是线性化可行方向集的子集
\end{itemize}
\end{notebox}

\subsection{二阶必要条件}

\begin{theorem}[二阶必要条件]
设 $f, g_i, h_j$ 二阶连续可微，$x^*$ 是局部极小点，LICQ成立，$(\lambda^*, \mu^*)$ 为对应KKT乘子。则：
\[
d^T \nabla_{xx}^2 \mathcal{L}(x^*, \lambda^*, \mu^*) d \geq 0, \quad \forall d \in \mathcal{C}(x^*)
\]
即Lagrange函数Hessian在临界锥上半正定。
\end{theorem}

\subsection{二阶充分条件}

\begin{theorem}[二阶充分条件]
设 $f, g_i, h_j$ 二阶连续可微，$x^*$ 满足KKT条件，对应乘子 $(\lambda^*, \mu^*)$。若：
\[
d^T \nabla_{xx}^2 \mathcal{L}(x^*, \lambda^*, \mu^*) d > 0, \quad \forall d \in \mathcal{C}(x^*) \setminus \{0\}
\]
则 $x^*$ 是严格局部极小点。
\end{theorem}

\begin{example}[二阶条件验证]
\[
\min \, x_1^2 + x_2^2, \quad \text{s.t. } x_1 + x_2 \geq 1
\]

\textbf{解：}

改写为 $g(x) = 1 - x_1 - x_2 \leq 0$。KKT解得 $x^* = (0.5, 0.5)$，$\lambda^* = 1$。

$\mathcal{L} = x_1^2 + x_2^2 + \lambda(1 - x_1 - x_2)$

$\nabla^2_{xx} \mathcal{L} = \begin{bmatrix} 2 & 0 \\ 0 & 2 \end{bmatrix}$

积极约束 $g$，$\lambda^* = 1 > 0$，临界锥：$\nabla g^T d = (-1,-1) \cdot (d_1, d_2) = 0$，即 $d_1 + d_2 = 0$。

对 $d = (t, -t)$：$d^T \nabla^2 \mathcal{L} d = 2t^2 + 2t^2 = 4t^2 > 0$（$t \neq 0$）

二阶充分条件满足，$x^* = (0.5, 0.5)$ 是严格局部极小点。
\end{example}

\subsection{例子：区分不同KKT点}

\begin{example}[多个KKT点的判断]
\[
\min \, x_1^3 + x_2^3, \quad \text{s.t. } x_1 + x_2 = 1
\]

\textbf{解：}

$\mathcal{L} = x_1^3 + x_2^3 + \mu(x_1 + x_2 - 1)$

KKT条件：
\begin{align*}
3x_1^2 + \mu &= 0 \\
3x_2^2 + \mu &= 0 \\
x_1 + x_2 &= 1
\end{align*}

由前两式：$x_1^2 = x_2^2 \Rightarrow x_1 = x_2$ 或 $x_1 = -x_2$

\textbf{Case 1}：$x_1 = x_2 = 0.5$，$\mu = -0.75$

$\nabla^2 \mathcal{L} = \begin{bmatrix} 6x_1 & 0 \\ 0 & 6x_2 \end{bmatrix} = \begin{bmatrix} 3 & 0 \\ 0 & 3 \end{bmatrix}$

临界锥：$d_1 + d_2 = 0$，即 $d = (t, -t)$

$d^T \nabla^2 \mathcal{L} d = 3t^2 + 3t^2 = 6t^2 > 0$ $\Rightarrow$ \textbf{严格局部极小}！

\textbf{Case 2}：$x_1 = -x_2$，结合 $x_1 + x_2 = 1$ 无解。

但若考虑 $x_1 = 0, x_2 = 1$（需验证是否为KKT点）：$\mu = 0$，但 $3(1)^2 + 0 \neq 0$，不是KKT点。
\end{example}

\begin{example}[KKT条件的几何验证]\label{ex:kkt-geometric}
\[
\min \, f(x) = (x_1 - 3)^2 + (x_2 - 2)^2, \quad \text{s.t. } x_1^2 + x_2^2 \leq 5,\ x_1 + 2x_2 \leq 4,\ x_1 \geq 0,\ x_2 \geq 0
\]

\textbf{解：}

\textbf{判断 $x' = (0, 0)$}：起作用约束为 $g_3 = -x_1$，$g_4 = -x_2$。$\nabla f(0,0) = (-6, -4)^T$，平稳性要求 $(-6, -4)^T + \lambda_3(-1, 0)^T + \lambda_4(0, -1)^T = 0$，得 $\lambda_3 = -6 < 0$，违反对偶可行性。故 $x'$ \textbf{不是 KKT 点}。

\textbf{判断 $x'' = (2, 1)$}：起作用约束为 $g_1 = x_1^2 + x_2^2 - 5$，$g_2 = x_1 + 2x_2 - 4$。$\nabla f(2,1) = (-2, -2)^T$，$\nabla g_1 = (4, 2)^T$，$\nabla g_2 = (1, 2)^T$。平稳性：
\[
(-2, -2)^T + \lambda_1(4, 2)^T + \lambda_2(1, 2)^T = 0 \Rightarrow \lambda_1 = \tfrac{1}{3},\ \lambda_2 = \tfrac{2}{3}
\]
均非负，互补松弛成立。故 $x''$ \textbf{是 KKT 点}。

\begin{center}
\begin{tikzpicture}[scale=1.0, >=Stealth]
  \draw[->] (-0.2, 0) -- (3.3, 0) node[right] {$x_1$};
  \draw[->] (0, -0.2) -- (0, 2.8) node[above] {$x_2$};
  % 约束曲线 x_1^2 + x_2^2 = 5（第一象限部分）
  \draw[domain=0:2.236, smooth, samples=100] plot(\x, {sqrt(5-\x*\x)}) node[right, scale=0.8] {$x_1^2+x_2^2=5$};
  % 约束直线 x_1 + 2x_2 = 4
  \draw[domain=0:3, smooth] plot(\x, {(4-\x)/2}) node[right, scale=0.8] {$x_1+2x_2=4$};
  % 可行域填充（简化为两约束交点附近区域）
  \fill[green!10] (0,0) -- (2,1) -- (2.236,0) -- cycle;
  % 标记点 x' = (0,0)
  \fill[red] (0, 0) circle (2.5pt) node[below left, scale=0.8] {$x'$};
  % 标记点 x'' = (2,1)
  \fill[green!50!black] (2, 1) circle (2.5pt) node[above right, scale=0.8] {$x''$};
  % 在 x'=(0,0) 处：∇f = 2(0-3, 0-2) = (-6, -4)，-∇f = (6, 4)
  \draw[->, thick, red] (0, 0) -- (0.9, 0.6) node[right, scale=0.7] {$-\nabla f(x')$};
  % 在 x''=(2,1) 处：∇f = 2(2-3, 1-2) = (-2, -2)，-∇f = (2, 2)
  \draw[->, thick, green!50!black] (2, 1) -- (2.6, 1.6) node[right, scale=0.7] {$-\nabla f(x'')$};
  % 起作用约束法向量（在 x'' 处）
  \draw[->, thick, blue!70!black] (2, 1) -- (2.4, 1.2) node[below, scale=0.65] {$\nabla g_1$};
  \draw[->, thick, orange] (2, 1) -- (2.2, 1.5) node[left, scale=0.65] {$\nabla g_2$};
\end{tikzpicture}
\end{center}

几何解释：$x'$ 处 $-\nabla f$ 指向可行域内部，不可能是最优点；$x''$ 位于两个约束边界交点，$-\nabla f(x'')$ 恰好可由起作用约束法向量的非负组合表示。
\end{example}

\begin{example}[二阶条件判定鞍点]
\[
\min \, x_1^2 - x_2^2, \quad \text{s.t. } x_1 + x_2 = 1
\]

\textbf{解：}

$\mathcal{L} = x_1^2 - x_2^2 + \mu(x_1 + x_2 - 1)$

KKT条件：$2x_1 + \mu = 0$，$-2x_2 + \mu = 0$，$x_1 + x_2 = 1$

解得：$x_1 = 0.5, x_2 = 0.5$，$\mu = -1$

$\nabla^2 \mathcal{L} = \begin{bmatrix} 2 & 0 \\ 0 & -2 \end{bmatrix}$

临界锥：$d_1 + d_2 = 0$，$d = (t, -t)$

$d^T \nabla^2 \mathcal{L} d = 2t^2 - 2t^2 = 0$

二阶条件无法判定。实际上这不是局部极小点（目标函数在约束线上无下界）。
\end{example}

% 第6章第4节：对偶理论
\section{对偶理论}\label{sec:duality}

\subsection{拉格朗日对偶问题}

\begin{definition}[对偶函数]
拉格朗日对偶函数定义为：
\[
q(\lambda, \mu) = \inf_{x \in \mathbb{R}^n} \mathcal{L}(x, \lambda, \mu) = \inf_{x \in \mathbb{R}^n} \left\{f(x) + \sum_{i=1}^m \lambda_i g_i(x) + \sum_{j=1}^l \mu_j h_j(x)\right\}
\]
\end{definition}

\begin{definition}[拉格朗日对偶问题]
\begin{align}
\max \quad & q(\lambda, \mu) \label{eq:dual} \\
\text{s.t.} \quad & \lambda \geq 0 \notag
\end{align}
\end{definition}

\begin{notebox}[对偶函数的性质]
\begin{itemize}
    \item $q(\lambda, \mu)$ 是凹函数（即使原问题非凸）
    \item 对偶问题是凸优化问题（最大化凹函数）
    \item 若 $\inf$ 在某点达到，则 $q(\lambda, \mu) = \mathcal{L}(x^*, \lambda, \mu)$
    \item 若 $\inf$ 未达到，$q(\lambda, \mu)$ 可取 $-\infty$
\end{itemize}
\end{notebox}

\subsection{弱对偶定理}

\begin{theorem}[弱对偶定理]
设 $\bar{x}$ 是原问题可行解，$(\lambda, \mu)$ 是对偶问题可行解（$\lambda \geq 0$），则：
\[
q(\lambda, \mu) \leq f(\bar{x})
\]
特别地，若 $p^*$ 和 $d^*$ 分别是原问题和对偶问题的最优值，则：
\[
d^* \leq p^*
\]
\end{theorem}

\textbf{证明}：对任意可行 $\bar{x}$ 和 $\lambda \geq 0$：
\begin{align*}
q(\lambda, \mu) &= \inf_x \mathcal{L}(x, \lambda, \mu) \leq \mathcal{L}(\bar{x}, \lambda, \mu) \\
&= f(\bar{x}) + \underbrace{\sum_i \lambda_i g_i(\bar{x})}_{\leq 0} + \underbrace{\sum_j \mu_j h_j(\bar{x})}_{= 0} \leq f(\bar{x})
\end{align*}

\begin{definition}[对偶间隙]
对偶间隙定义为 $p^* - d^* \geq 0$。若 $p^* = d^*$，称\textbf{强对偶}成立。
\end{definition}

\subsection{强对偶定理}

\begin{theorem}[强对偶定理]
对于凸优化问题，若Slater条件成立（存在严格内点），则强对偶成立，即 $p^* = d^*$。

若 $p^*$ 有限，则对偶最优解 $(\lambda^*, \mu^*)$ 可达。
\end{theorem}

\begin{example}[强对偶成立]
\[
\min \, x_1^2 + x_2^2, \quad \text{s.t. } x_1 + x_2 \geq 1
\]

\textbf{解：}

改写为 $g(x) = 1 - x_1 - x_2 \leq 0$。

$\mathcal{L} = x_1^2 + x_2^2 + \lambda(1 - x_1 - x_2)$

对 $\lambda \geq 0$ 求 $\inf$：
\[
\frac{\partial \mathcal{L}}{\partial x_1} = 2x_1 - \lambda = 0 \Rightarrow x_1 = \lambda/2
\]
同理 $x_2 = \lambda/2$。

$q(\lambda) = \frac{\lambda^2}{4} + \frac{\lambda^2}{4} + \lambda(1 - \lambda) = \lambda - \frac{\lambda^2}{2}$

对偶问题：$\max_{\lambda \geq 0} \, \lambda - \frac{\lambda^2}{2}$

$q'(\lambda) = 1 - \lambda = 0 \Rightarrow \lambda^* = 1$

$d^* = q(1) = 1 - 0.5 = 0.5 = p^*$。强对偶成立。
\end{example}

\begin{example}[存在对偶间隙]
\[
\min \, -x, \quad \text{s.t. } x^3 \leq 0, \; x \in \mathbb{R}
\]

\textbf{解：}

等价于 $x \leq 0$，最优解 $x^* = 0$，$p^* = 0$。

$q(\lambda) = \inf_x \{-x + \lambda x^3\}$

当 $\lambda > 0$：令 $x \to -\infty$，$\lambda x^3 - x \to -\infty$，$q(\lambda) = -\infty$

当 $\lambda = 0$：$q(0) = \inf_x (-x) = -\infty$

$d^* = \sup_{\lambda \geq 0} q(\lambda) = -\infty < p^* = 0$，存在对偶间隙！

原因：非凸约束（$x^3 \leq 0$ 等价于 $x \leq 0$，但写成 $x^3$ 非凸），Slater条件不满足。
\end{example}

\subsection{鞍点定理}

\begin{theorem}[鞍点定理]
$(x^*, \lambda^*, \mu^*)$ 是 $\mathcal{L}$ 的鞍点当且仅当：
\begin{enumerate}
    \item $x^*$ 是原问题最优解
    \item $(\lambda^*, \mu^*)$ 是对偶问题最优解
    \item $p^* = d^*$（无对偶间隙）
\end{enumerate}
\end{theorem}

\subsection{扰动函数与灵敏度分析}

\begin{definition}[扰动函数]
定义 $v(u) = \inf\{f(x) : g(x) \leq u, \, h(x) = 0\}$，其中 $u$ 是约束右端项的扰动。
\end{definition}

\begin{theorem}[灵敏度分析]
若强对偶成立，则对偶最优乘子 $\lambda^*$ 满足：
\[
\lambda_i^* = -\frac{\partial v(0)}{\partial u_i}
\]
即 $\lambda_i^*$ 表示约束右端项增加1单位时，最优值的减小率（影子价格）。
\end{theorem}

\begin{example}[灵敏度分析]
\[
\min \, x_1^2 + x_2^2, \quad \text{s.t. } x_1 + x_2 \geq 1
\]

\textbf{解：}

扰动约束为 $x_1 + x_2 \geq 1 + u$。$v(u) = \frac{(1+u)^2}{2}$（对称解）

$\lambda^* = 1$（前面已求得）

$v'(0) = (1+u)|_{u=0} = 1 = -\lambda^*$ (验证一致)
\end{example}

% 第6章第5节：惩罚函数法
\section{惩罚函数法}\label{sec:ch6-penalty}

\subsection{二次惩罚函数法}

\begin{conceptbox}[基本思想]
将约束 violation 以惩罚项形式加入目标函数，通过增大惩罚参数 $\rho \to \infty$，使无约束问题的解逼近原约束问题的解。
\end{conceptbox}

\begin{definition}[二次惩罚函数]
对于问题 \eqref{eq:nlp}，定义惩罚函数：
\[
P(x, \rho) = f(x) + \frac{\rho}{2}\left[\sum_{i=1}^m \max(0, g_i(x))^2 + \sum_{j=1}^l h_j(x)^2\right]
\]
其中 $\rho > 0$ 为惩罚参数。
\end{definition}

\begin{figure}[H]
\centering
\begin{tikzpicture}[scale=0.72]
    % 坐标轴
    \draw[->,thick] (-0.3,0) -- (7.5,0) node[right] {$x$};
    \draw[->,thick] (0,-0.3) -- (0,5.5) node[above] {$f(x)$};
    % 原目标函数
    \draw[very thick,blue]
        (0.5,4.5) .. controls (1.5,2) and (2.5,0.8) .. (3.5,0.5)
        .. controls (4.5,0.8) and (5.5,2) .. (6.5,4);
    \node[blue,scale=0.8] at (6.7,4.2) {$f$};
    % 约束边界
    \draw[red,thick] (4,0) -- (4,5) node[above,scale=0.75] {\small $g(x)\!=\!0$};
    \fill[red!10] (4,0) rectangle (7.5,5);
    \node[red,scale=0.7] at (5.5,4.5) {\small 不可行域};
    % 惩罚函数（不同 rho）
    \draw[orange,thick]
        (0.5,4.5) .. controls (1.5,2) and (2.5,0.8) .. (3.2,0.6)
        .. controls (3.5,0.7) and (3.8,1.5) .. (4,3);
    \node[orange,scale=0.7,right] at (4,3) {\small $\rho$ 小};
    \draw[green!50!black,thick]
        (0.5,4.5) .. controls (1.5,2) and (2.5,0.8) .. (3.4,0.55)
        .. controls (3.6,0.6) and (3.9,2) .. (4,5);
    \node[green!50!black,scale=0.7,right] at (4,5) {\small $\rho$ 大};
    % 最优解
    \fill[black] (3.5,0.5) circle (2.5pt) node[below,scale=0.75] {\small $x^*$};
    \draw[dotted] (3.5,0) -- (3.5,0.5);
    % 惩罚解
    \fill[orange] (3.2,0.6) circle (2pt);
    \fill[green!50!black] (3.4,0.55) circle (2pt);
    \draw[->,thick,gray] (3.2,0.6) -- (3.45,0.55) -- (3.5,0.5);
    \node[gray,scale=0.65,above] at (3.35,0.7) {\small $\rho\!\to\!\infty$};
\end{tikzpicture}
\caption{二次惩罚函数法：增大 $\rho$ 使惩罚函数的极小点逼近约束最优解 $x^*$}
\label{fig:penalty}
\end{figure}

\begin{algobox}[二次惩罚法算法]
\textbf{初始化}：选取 $x^{(0)}$，$\rho_0 > 0$，放大因子 $\beta > 1$，容差 $\varepsilon > 0$

\textbf{For} $k = 0, 1, 2, \ldots$：
\begin{enumerate}
    \item 以 $x^{(k)}$ 为初始点，求解 $\min_x P(x, \rho_k)$，得 $x^{(k+1)}$
    \item 若约束 violation $< \varepsilon$，停止
    \item 更新 $\rho_{k+1} = \beta \rho_k$
\end{enumerate}
\end{algobox}

\begin{theorem}[二次惩罚法的收敛性]
设 $f, g_i, h_j$ 连续，$\rho_k \to \infty$。若 $x^{(k)}$ 是 $P(x, \rho_k)$ 的全局极小点，则 $\{x^{(k)}\}$ 的任何聚点都是原问题的全局最优解。
\end{theorem}

\begin{notebox}[二次惩罚法的缺点]
\begin{itemize}
    \item 当 $\rho$ 很大时，Hessian条件数恶化，数值困难
    \item 需要求解一系列无约束问题，计算量大
    \item 对非凸问题可能收敛到局部极小
\end{itemize}
\end{notebox}

\begin{example}[二次惩罚法——迭代两次]
\[
\min \, x^2, \quad \text{s.t. } x \geq 1
\]

\textbf{解：}

惩罚函数：$P(x, \rho) = x^2 + \frac{\rho}{2}\max(0, 1-x)^2$

\textbf{第1次迭代}（$\rho_0 = 1$，初值 $x^{(0)} = 0$）：

$P(x, 1) = x^2 + 0.5(1-x)^2$（假设 $x < 1$）

$\frac{dP}{dx} = 2x - (1-x) = 3x - 1 = 0 \Rightarrow x^{(1)} = \frac{1}{3} \approx 0.333$

约束 violation：$1 - 1/3 = 2/3 > \varepsilon$，继续。

\textbf{第2次迭代}（$\rho_1 = 10$）：

$P(x, 10) = x^2 + 5(1-x)^2$

$\frac{dP}{dx} = 2x - 10(1-x) = 12x - 10 = 0 \Rightarrow x^{(2)} = \frac{10}{12} = \frac{5}{6} \approx 0.833$

可见随着 $\rho$ 增大，$x^{(k)} \to 1$（约束最优解）。
\end{example}

\subsection{精确惩罚函数法}

\begin{conceptbox}[精确惩罚]
二次惩罚法需要 $\rho \to \infty$。精确惩罚函数法通过$\ell_1$惩罚，对有限 $\rho$ 即可得到精确解。
\end{conceptbox}

\begin{definition}[$\ell_1$精确惩罚函数]
\[
\Phi(x, \rho) = f(x) + \rho\left[\sum_{i=1}^m \max(0, g_i(x)) + \sum_{j=1}^l |h_j(x)|\right]
\]
\end{definition}

\begin{theorem}[精确性]
设 $x^*$ 是原问题的局部最优解，LICQ和严格互补条件成立。则存在 $\bar{\rho} > 0$，使得对所有 $\rho \geq \bar{\rho}$，$x^*$ 也是 $\Phi(x, \rho)$ 的局部极小点。
\end{theorem}

\begin{example}[精确惩罚法——迭代两次]
\[
\min \, x^2 + 2x, \quad \text{s.t. } x \geq 1
\]

\textbf{解：}

精确惩罚：$\Phi(x, \rho) = x^2 + 2x + \rho \max(0, 1-x)$

\textbf{Case A}（$x \geq 1$）：$\Phi = x^2 + 2x$，$\Phi' = 2x + 2 > 0$（$x \geq 1$），最小在 $x = 1$。

\textbf{Case B}（$x < 1$）：$\Phi = x^2 + 2x + \rho(1-x) = x^2 + (2-\rho)x + \rho$

$\Phi' = 2x + 2 - \rho = 0 \Rightarrow x = \frac{\rho - 2}{2}$

\textbf{第1次迭代}（$\rho = 3$）：$x = 0.5 < 1$，Case B。$\Phi(0.5) = 0.25 + 1 + 1.5 = 2.75$。

与Case A的边界值 $\Phi(1) = 3$ 比较，最小值在 $x = 0.5$（不精确）。

\textbf{第2次迭代}（$\rho = 5$）：$x = 1.5 \geq 1$，转入Case A。

在Case A中，$\Phi' = 2x + 2 > 0$（$x \geq 1$），故 $x^* = 1$ 是最小点。

$\Phi(1) = 3$，验证了 $x = 1$ 是精确解。
\end{example}

\subsection{增广拉格朗日法}

增广拉格朗日法结合了拉格朗日函数和惩罚函数的优点，对有限惩罚参数即可收敛。

\begin{definition}[增广拉格朗日函数]
对于等式约束问题：
\[
\mathcal{L}_A(x, \mu, \rho) = f(x) + \mu^T h(x) + \frac{\rho}{2}\|h(x)\|^2
\]
\end{definition}

\begin{algobox}[增广拉格朗日法]
\textbf{初始化}：$x^{(0)}$，$\mu^{(0)}$，$\rho_0 > 0$

\textbf{For} $k = 0, 1, 2, \ldots$：
\begin{enumerate}
    \item $x^{(k+1)} = \arg\min_x \mathcal{L}_A(x, \mu^{(k)}, \rho_k)$
    \item $\mu^{(k+1)} = \mu^{(k)} + \rho_k h(x^{(k+1)})$
    \item 若 $\|h(x^{(k+1)})\| < \varepsilon$，停止
    \item 可选：$\rho_{k+1} = \min(\beta \rho_k, \rho_{\max})$
\end{enumerate}
\end{algobox}

\begin{notebox}[乘子更新原理]
由KKT条件 $\nabla f + \mu^T \nabla h = 0$ 和增广拉格朗日函数的平稳性：
\[
\nabla f + (\mu^{(k)})^T \nabla h + \rho_k h^T \nabla h = 0
\]
对比得：$\mu^{(k+1)} = \mu^{(k)} + \rho_k h(x^{(k+1)})$ 是对真实乘子的估计。
\end{notebox}

\begin{example}[增广拉格朗日法——迭代两次]
\[
\min \, x^2, \quad \text{s.t. } x = 2
\]

\textbf{解：}

$\mathcal{L}_A = x^2 + \mu(x-2) + \frac{\rho}{2}(x-2)^2$

$\frac{\partial \mathcal{L}_A}{\partial x} = 2x + \mu + \rho(x-2) = 0 \Rightarrow x = \frac{2\rho - \mu}{2 + \rho}$

\textbf{第1次迭代}（$\mu^{(0)} = 0$，$\rho_0 = 1$）：

$x^{(1)} = \frac{2 - 0}{3} = \frac{2}{3} \approx 0.667$

$\mu^{(1)} = 0 + 1 \cdot (2/3 - 2) = -4/3 \approx -1.333$

\textbf{第2次迭代}（$\rho_1 = 2$）：

$x^{(2)} = \frac{4 - (-4/3)}{4} = \frac{16/3}{4} = \frac{4}{3} \approx 1.333$

$\mu^{(2)} = -4/3 + 2(4/3 - 2) = -4/3 - 4/3 = -8/3 \approx -2.667$

继续迭代，$x^{(k)} \to 2$，$\mu^{(k)} \to -4$（真实乘子：$\nabla f = 2x = 4$，$\mu = -4$）。
\end{example}

\subsection{Python代码实现}\label{sec:penalty-code}

\begin{codebox}[二次惩罚法]
\begin{lstlisting}
import numpy as np
from scipy.optimize import minimize

def quadratic_penalty(f, grad_f, g_funcs, h_funcs, x0,
                      rho0=1.0, beta=10.0, eps=1e-6,
                      max_iter=50):
    """
    二次惩罚函数法
    
    Parameters:
        f: 目标函数
        grad_f: 目标函数梯度
        g_funcs: 不等式约束函数列表 (g_i <= 0)
        h_funcs: 等式约束函数列表
        x0: 初始点
        rho0: 初始惩罚参数
        beta: 放大因子
        eps: 约束容差
    
    Returns:
        x_opt: 最优解
        history: 迭代历史
    """
    x = x0.copy()
    rho = rho0
    history = []
    
    for k in range(max_iter):
        # 构造惩罚函数
        def P(x_val):
            penalty = 0.0
            for g in g_funcs:
                gv = g(x_val)
                if gv > 0:
                    penalty += gv**2
            for h in h_funcs:
                penalty += h(x_val)**2
            return f(x_val) + 0.5 * rho * penalty
        
        # 求解无约束问题
        res = minimize(P, x, method='BFGS')
        x = res.x
        
        # 计算约束违反
        violation = 0.0
        for g in g_funcs:
            violation += max(0, g(x))**2
        for h in h_funcs:
            violation += h(x)**2
        violation = np.sqrt(violation)
        
        history.append((x.copy(), f(x), violation, rho))
        
        if violation < eps:
            break
        
        rho *= beta
    
    return x, history

# 示例: min x^2, s.t. x >= 1
f = lambda x: x[0]**2
g = [lambda x: 1 - x[0]]  # 1 - x <= 0, i.e., x >= 1

x0 = np.array([0.0])
x_opt, hist = quadratic_penalty(f, None, g, [], x0,
                                rho0=1.0, beta=10.0)
for i, (x, fx, v, rho) in enumerate(hist):
    print(f"Iter {i}: x={x[0]:.4f}, f={fx:.4f}, "
          f"viol={v:.4f}, rho={rho:.1f}")
\end{lstlisting}
\end{codebox}

\begin{codebox}[增广拉格朗日法]
\begin{lstlisting}
def augmented_lagrangian(f, h_funcs, x0, mu0, rho0=1.0,
                          beta=2.0, eps=1e-6, max_iter=50):
    """
    增广拉格朗日法（等式约束）
    
    Parameters:
        f: 目标函数
        h_funcs: 等式约束函数列表
        x0: 初始点
        mu0: 初始乘子
        rho0: 初始惩罚参数
        beta: 放大因子
    
    Returns:
        x_opt: 最优解
        mu_opt: 最优乘子
        history: 迭代历史
    """
    x = np.array(x0, dtype=float)
    mu = np.array(mu0, dtype=float)
    rho = rho0
    history = []
    
    for k in range(max_iter):
        # 构造增广拉格朗日函数
        def L_A(x_val):
            h_vals = np.array([h(x_val) for h in h_funcs])
            return f(x_val) + np.dot(mu, h_vals) + \
                   0.5 * rho * np.sum(h_vals**2)
        
        res = minimize(L_A, x, method='BFGS')
        x = res.x
        
        h_vals = np.array([h(x) for h in h_funcs])
        h_norm = np.linalg.norm(h_vals)
        
        history.append((x.copy(), mu.copy(), rho, h_norm))
        
        if h_norm < eps:
            break
        
        mu = mu + rho * h_vals
        rho = min(beta * rho, 1e6)
    
    return x, mu, history

# 示例: min x^2, s.t. x = 2
f = lambda x: x[0]**2
h = [lambda x: x[0] - 2]

x0 = np.array([0.0])
mu0 = np.array([0.0])
x_opt, mu_opt, hist = augmented_lagrangian(f, h, x0, mu0)
for i, (x, mu, rho, h_n) in enumerate(hist):
    print(f"Iter {i}: x={x[0]:.4f}, mu={mu[0]:.4f}, "
          f"rho={rho:.1f}, ||h||={h_n:.4f}")
\end{lstlisting}
\end{codebox}

% 第6章第6节：障碍函数法与内点法
\section{障碍函数法与内点法}\label{sec:interior}

\subsection{对数障碍函数法}

\begin{conceptbox}[内点法思想]
内点法通过在可行域内部迭代，用障碍函数阻止迭代点靠近边界。只适用于不等式约束。
\end{conceptbox}

\begin{definition}[对数障碍函数]
对于不等式约束 $g_i(x) \leq 0$（即 $-g_i(x) \geq 0$），定义对数障碍函数：
\[
B(x, \mu) = f(x) - \mu \sum_{i=1}^m \ln(-g_i(x))
\]
其中 $\mu > 0$ 为障碍参数。当 $x$ 接近边界时 $B \to +\infty$。
\end{definition}

\begin{algobox}[对数障碍法]
\textbf{初始化}：严格可行点 $x^{(0)}$（$g_i(x^{(0)}) < 0$），$\mu_0 > 0$，缩减因子 $\sigma \in (0,1)$

\textbf{For} $k = 0, 1, 2, \ldots$：
\begin{enumerate}
    \item 以 $x^{(k)}$ 为初始点，求解 $\min_x B(x, \mu_k)$，得 $x^{(k+1)}$
    \item 若 $\mu_k < \varepsilon$，停止
    \item $\mu_{k+1} = \sigma \mu_k$
\end{enumerate}
\end{algobox}

\begin{theorem}[收敛性]
设 $f, g_i$ 连续，$\mu_k \to 0$。若 $x^{(k)}$ 是 $B(x, \mu_k)$ 的全局极小点，则 $\{x^{(k)}\}$ 的任何聚点都是原问题的全局最优解。
\end{theorem}

\begin{notebox}[障碍函数法的特点]
\begin{itemize}
    \item 只需要一个初始严格可行点
    \item 迭代点始终保持在可行域内部
    \item 当 $\mu \to 0$ 时，条件数恶化（类似惩罚法）
    \item 适用于凸优化问题效果良好
\end{itemize}
\end{notebox}

\begin{example}[对数障碍法——迭代两次]
\[
\min \, x, \quad \text{s.t. } x \geq 0, \; x \leq 1
\]

\textbf{解：}

改写：$g_1(x) = -x \leq 0$，$g_2(x) = x - 1 \leq 0$

障碍函数：$B(x, \mu) = x - \mu\ln(x) - \mu\ln(1-x)$（要求 $0 < x < 1$）

$\frac{dB}{dx} = 1 - \frac{\mu}{x} + \frac{\mu}{1-x} = 0$

$\Rightarrow x(1-x) - \mu(1-x) + \mu x = 0$

$\Rightarrow x - x^2 - \mu + \mu x + \mu x = 0$

$\Rightarrow x^2 - (1+2\mu)x + \mu = 0$

\textbf{第1次迭代}（$\mu_0 = 0.1$）：

$x^2 - 1.2x + 0.1 = 0$

$x = \frac{1.2 - \sqrt{1.44 - 0.4}}{2} = \frac{1.2 - 1.02}{2} \approx 0.09$
（取在 $(0,1)$ 内的根）

实际解：$x^{(1)} = \frac{1.2 - \sqrt{1.04}}{2} \approx \frac{1.2 - 1.0199}{2} \approx 0.090$

\textbf{第2次迭代}（$\mu_1 = 0.01$）：

$x^2 - 1.02x + 0.01 = 0$

$x = \frac{1.02 - \sqrt{1.0404 - 0.04}}{2} = \frac{1.02 - 1.0}{2} \approx 0.010$

可见 $x^{(k)} \to 0$（最优解）。
\end{example}

\subsection{中心路径}

\begin{definition}[中心路径]
对于凸优化问题，当 $\mu$ 从 $\infty$ 变到 $0$ 时，$B(x, \mu)$ 的极小点 $x(\mu)$ 形成一条轨迹，称为\textbf{中心路径}（Central Path）。
\end{definition}

\begin{theorem}[中心路径的性质]
\begin{enumerate}
    \item $x(\mu)$ 在可行域严格内部
    \item $\mu \to 0$ 时，$x(\mu) \to x^*$（最优解）
    \item $\mu \to \infty$ 时，$x(\mu) \to$ 解析中心
\end{enumerate}
\end{theorem}

\subsection{原始-对偶内点法}

\begin{conceptbox}[原始-对偶方法]
同时更新原始变量 $x$ 和对偶变量 $\lambda$，沿着牛顿方向逼近KKT条件。
\end{conceptbox}

考虑问题：$\min f(x)$ s.t. $g_i(x) \leq 0$。

KKT条件改写（引入松弛变量 $s_i = -g_i(x) \geq 0$）：
\begin{align*}
\nabla f(x) + \sum_i \lambda_i \nabla g_i(x) &= 0 \\
\lambda_i s_i &= \mu \quad \text{(扰动互补)} \\
s_i + g_i(x) &= 0 \\
\lambda_i, s_i &\geq 0
\end{align*}

当 $\mu \to 0$ 时，逼近精确KKT条件。

\begin{algobox}[原始-对偶内点法]
\textbf{初始化}：严格可行点 $x^{(0)}$，$\lambda^{(0)} > 0$，$s^{(0)} > 0$，$\mu_0 > 0$

\textbf{For} $k = 0, 1, 2, \ldots$：
\begin{enumerate}
    \item 求解牛顿方程组得到方向 $(\Delta x, \Delta \lambda, \Delta s)$
    \item 线搜索确定步长 $\alpha$，保持 $\lambda, s > 0$
    \item 更新：$(x^{(k+1)}, \lambda^{(k+1)}, s^{(k+1)}) = (x^{(k)}, \lambda^{(k)}, s^{(k)}) + \alpha(\Delta x, \Delta \lambda, \Delta s)$
    \item $\mu_{k+1} = \sigma \mu_k$（$\sigma \in (0,1)$）
    \item 若 duality gap $< \varepsilon$，停止
\end{enumerate}
\end{algobox}

\subsection{线性规划的内点法}

\begin{example}[LP内点法——迭代两次]
\[
\min \, -x_1 - 2x_2, \quad \text{s.t. } x_1 + x_2 \leq 2, \; x_1, x_2 \geq 0
\]

\textbf{解：}

改写：$\min c^T x$ s.t. $Ax + s = b$，$x, s \geq 0$

其中 $c = (-1, -2)^T$，$A = [1, 1]$，$b = 2$。

障碍函数：$B = c^T x - \mu(\ln x_1 + \ln x_2 + \ln s_1)$

约束：$x_1 + x_2 + s_1 = 2$

消去 $s_1 = 2 - x_1 - x_2$，对 $x_1, x_2$ 求极小：

$\frac{\partial B}{\partial x_1} = -1 - \frac{\mu}{x_1} + \frac{\mu}{s_1} = 0$

$\frac{\partial B}{\partial x_2} = -2 - \frac{\mu}{x_2} + \frac{\mu}{s_1} = 0$

由前两式相减：$1 - \frac{\mu}{x_1} + \frac{\mu}{x_2} = 0$

\textbf{第1次迭代}（$\mu_0 = 0.5$）：

设 $x_1 = x_2 = t$，则 $1 = 0$ 矛盾。尝试数值求解：

设 $x_1 \approx 0.4, x_2 \approx 0.8, s_1 = 0.8$

验证：$-1 - 1.25 + 0.625 = -1.625 \neq 0$，需迭代调整。

经数值求解（如牛顿法），$x^{(1)} \approx (0.35, 0.90)$，$s_1^{(1)} \approx 0.75$

\textbf{第2次迭代}（$\mu_1 = 0.1$）：

类似求解，$x^{(2)} \approx (0.10, 1.35)$，更接近边界 $x_1 = 0$。

最优解为 $(0, 2)$，最优值 $-4$。
\end{example}

\subsection{Python代码实现}\label{sec:interior-code}

\begin{codebox}[对数障碍法]
\begin{lstlisting}
import numpy as np
from scipy.optimize import minimize

def log_barrier_method(f, g_funcs, x0, mu0=1.0,
                       sigma=0.1, eps=1e-6, max_iter=50):
    """
    对数障碍函数法
    
    Parameters:
        f: 目标函数
        g_funcs: 不等式约束 g_i(x) <= 0 列表
        x0: 严格可行初始点
        mu0: 初始障碍参数
        sigma: 缩减因子
        eps: 精度
    
    Returns:
        x_opt: 最优解
        history: 迭代历史
    """
    x = np.array(x0, dtype=float)
    mu = mu0
    history = []
    
    for k in range(max_iter):
        def B(x_val):
            barrier = 0.0
            for g in g_funcs:
                gv = g(x_val)
                if gv >= 0:
                    return 1e10
                barrier -= np.log(-gv)
            return f(x_val) + mu * barrier
        
        res = minimize(B, x, method='BFGS',
                      options={'gtol': 1e-8})
        x = res.x
        
        history.append((x.copy(), mu, f(x)))
        
        if mu < eps:
            break
        
        mu *= sigma
    
    return x, history

# 示例: min x, s.t. x >= 0, x <= 1
# Rewrite: g1(x) = -x <= 0, g2(x) = x - 1 <= 0
f = lambda x: x[0]
g = [lambda x: -x[0], lambda x: x[0] - 1]

x0 = np.array([0.5])
x_opt, hist = log_barrier_method(f, g, x0, mu0=0.1, sigma=0.1)
for i, (x, mu, fx) in enumerate(hist):
    print(f"Iter {i}: x={x[0]:.6f}, mu={mu:.4f}, f={fx:.6f}")
\end{lstlisting}
\end{codebox}

\begin{codebox}[使用CVXOPT的原始-对偶内点法]
\begin{lstlisting}
import numpy as np

def primal_dual_lp(c, A, b, x0, s0, lam0,
                   mu0=1.0, sigma=0.1, eps=1e-8,
                   max_iter=100):
    """
    原始-对偶内点法求解标准形式LP:
    min c^T x, s.t. Ax = b, x >= 0
    
    Parameters:
        c: 目标函数系数
        A, b: 等式约束
        x0, s0, lam0: 初始点 (x > 0, s > 0)
        mu0: 初始中心参数
        sigma: 缩减因子
    
    Returns:
        x_opt, history
    """
    n = len(c)
    m = len(b)
    
    x = np.array(x0, dtype=float)
    s = np.array(s0, dtype=float)
    lam = np.array(lam0, dtype=float)
    mu = mu0
    history = []
    
    for k in range(max_iter):
        # 残差
        r_b = A @ x - b
        r_c = A.T @ lam + s - c
        r_mu = x * s - mu
        
        gap = np.dot(x, s) / n
        history.append((x.copy(), gap, mu))
        
        if gap < eps and np.linalg.norm(r_b) < eps:
            break
        
        # 牛顿方程
        X_inv = 1.0 / x
        S = s
        
        # 简化: 解正规方程
        D = x / s
        M = A @ np.diag(D) @ A.T
        
        rhs = -r_b - A @ (D * (r_c - r_mu / x))
        d_lam = np.linalg.solve(M, rhs)
        d_s = -r_c - A.T @ d_lam
        d_x = -(r_mu + x * d_s) / s
        
        # 最大步长保持正性
        alpha_x = min(1.0, 0.99 * min(
            [-xi/dxi for xi, dxi in zip(x, d_x) if dxi < 0],
            default=1.0))
        alpha_s = min(1.0, 0.99 * min(
            [-si/dsi for si, dsi in zip(s, d_s) if dsi < 0],
            default=1.0))
        
        x += alpha_x * d_x
        s += alpha_s * d_s
        lam += alpha_s * d_lam
        mu = sigma * gap
    
    return x, history

# 示例: min -x1 - 2*x2, s.t. x1 + x2 + x3 = 2, x >= 0
c = np.array([-1.0, -2.0, 0.0])
A = np.array([[1.0, 1.0, 1.0]])
b = np.array([2.0])

x0 = np.array([0.5, 0.5, 1.0])
s0 = np.array([1.0, 1.0, 1.0])
lam0 = np.array([0.0])

x_opt, hist = primal_dual_lp(c, A, b, x0, s0, lam0)
for i, (x, gap, mu) in enumerate(hist[:5]):
    print(f"Iter {i}: x=({x[0]:.4f},{x[1]:.4f}), "
          f"gap={gap:.6f}, mu={mu:.6f}")
\end{lstlisting}
\end{codebox}

% 第6章第7节：线性规划对偶与共轭函数
\section{线性规划对偶与Fenchel对偶}\label{sec:lp-conjugate}

\subsection{线性规划的对偶}

标准形式的线性规划：
\begin{align}
\min \quad & c^T x \label{eq:lp-primal} \\
\text{s.t.} \quad & Ax = b \notag \\
& x \geq 0 \notag
\end{align}

\begin{theorem}[LP对偶]
线性规划 \eqref{eq:lp-primal} 的对偶问题为：
\begin{align}
\max \quad & b^T y \label{eq:lp-dual} \\
\text{s.t.} \quad & A^T y \leq c \notag
\end{align}
\end{theorem}

\textbf{推导}：

拉格朗日函数：$\mathcal{L}(x, y, s) = c^T x - y^T(Ax - b) - s^T x = (c - A^T y - s)^T x + b^T y$

对偶函数：$q(y, s) = \inf_{x} \mathcal{L}(x, y, s) = \begin{cases} b^T y & \text{if } c - A^T y - s = 0 \\ -\infty & \text{otherwise} \end{cases}$

对偶问题：$\max_{y, s \geq 0} \{b^T y : A^T y + s = c\}$，即 $\max_y \{b^T y : A^T y \leq c\}$。

\begin{example}[LP对偶]
原问题：
\[
\min \, 3x_1 + 2x_2, \quad \text{s.t. } x_1 + x_2 = 2, \; x_1, x_2 \geq 0
\]

\textbf{解：}

对偶问题：
\[
\max \, 2y, \quad \text{s.t. } y \leq 3, \; y \leq 2
\]

对偶最优：$y^* = 2$，$d^* = 4$。

原问题最优：$x^* = (0, 2)$，$p^* = 4 = d^*$。
\end{example}

\subsection{LP对偶的基本定理}

\begin{theorem}[LP对偶定理]
对于线性规划原问题和对偶问题：
\begin{enumerate}
    \item 若原问题和对偶问题都有可行解，则都有最优解，且 $p^* = d^*$
    \item 若一方无界，则另一方不可行
    \item 若一方不可行，另一方可能无界或不可行
\end{enumerate}
\end{theorem}

\begin{notebox}[互补松弛条件（LP）]
对于LP，KKT互补松弛条件简化为：
\[
x_i \cdot (c_i - a_i^T y) = 0, \quad i = 1, \ldots, n
\]
即：若 $x_i > 0$，则对应的对偶约束取等号（$a_i^T y = c_i$）；若 $a_i^T y < c_i$，则 $x_i = 0$。
\end{notebox}

\subsection{共轭函数}

\begin{definition}[Fenchel共轭函数]
函数 $f: \mathbb{R}^n \to \mathbb{R} \cup \{+\infty\}$ 的共轭函数定义为：
\[
f^*(y) = \sup_{x \in \mathbb{R}^n} \{y^T x - f(x)\}
\]
\end{definition}

\begin{theorem}[共轭函数的性质]
\begin{enumerate}
    \item $f^*$ 是凸函数（闭的下半连续凸包络）
    \item 若 $f$ 是闭的真凸函数，则 $f^{**} = f$
    \item （Fenchel不等式）$f(x) + f^*(y) \geq x^T y$
    \item 等号成立 $\Leftrightarrow y \in \partial f(x)$
\end{enumerate}
\end{theorem}

\begin{example}[常见共轭函数]
\textbf{解：}

\begin{enumerate}
    \item $f(x) = \frac{1}{2}x^2$：$f^*(y) = \frac{1}{2}y^2$
    \item $f(x) = |x|$：$f^*(y) = \begin{cases} 0 & |y| \leq 1 \\ +\infty & |y| > 1 \end{cases}$
    \item $f(x) = \begin{cases} 0 & x \in C \\ +\infty & x \notin C \end{cases}$（示性函数）：$f^*(y) = \sup_{x \in C} y^T x$（支撑函数）
    \item $f(x) = \frac{1}{2}x^T Q x$（$Q \succ 0$）：$f^*(y) = \frac{1}{2}y^T Q^{-1} y$
\end{enumerate}
\end{example}

\subsection{Fenchel对偶}

\begin{theorem}[Fenchel对偶]
考虑问题：
\[
\min \, f(x) + g(Ax)
\]
其Fenchel对偶为：
\[
\max \, -f^*(A^T y) - g^*(-y)
\]
若 $f$ 和 $g$ 为闭真凸函数，且存在相对内点条件，则强对偶成立。
\end{theorem}

\begin{example}[Fenchel对偶]
$\min \, \frac{1}{2}\|x\|^2 + \|Ax - b\|_1$

\textbf{解：}

令 $f(x) = \frac{1}{2}\|x\|^2$，$g(z) = \|z - b\|_1$

$f^*(y) = \frac{1}{2}\|y\|^2$，$g^*(w) = \begin{cases} b^T w & \|w\|_\infty \leq 1 \\ +\infty & \text{otherwise} \end{cases}$

Fenchel对偶：
\[
\max_{\|y\|_\infty \leq 1} \left\{-\frac{1}{2}\|A^T y\|^2 + b^T y\right\}
\]
\end{example}

\subsection{正则化问题的对偶}

\begin{example}[L2正则化的对偶]
原问题（岭回归）：
\[
\min \, \frac{1}{2}\|Ax - b\|^2 + \frac{\lambda}{2}\|x\|^2
\]

\textbf{解：}

令 $z = Ax - b$，则问题变为：
\[
\min \, \frac{1}{2}\|z\|^2 + \frac{\lambda}{2}\|x\|^2 \quad \text{s.t. } z = Ax - b
\]

拉格朗日函数：$\mathcal{L} = \frac{1}{2}\|z\|^2 + \frac{\lambda}{2}\|x\|^2 + y^T(z - Ax + b)$

对 $z$ 求极小：$z + y = 0 \Rightarrow z = -y$

对 $x$ 求极小：$\lambda x - A^T y = 0 \Rightarrow x = \frac{1}{\lambda}A^T y$

代入得对偶：
\[
\max_y \left\{b^T y - \frac{1}{2}\|y\|^2 - \frac{1}{2\lambda}\|A^T y\|^2\right\}
\]
\end{example}

\begin{example}[L1正则化（LASSO）的对偶]
原问题：
\[
\min \, \frac{1}{2}\|Ax - b\|^2 + \lambda\|x\|_1
\]

\textbf{解：}

令 $f(x) = \frac{1}{2}\|Ax - b\|^2$，$g(x) = \lambda\|x\|_1$

$g^*(y) = \begin{cases} 0 & \|y\|_\infty \leq \lambda \\ +\infty & \text{otherwise} \end{cases}$

Fenchel对偶：
\[
\max_y \left\{-f^*(A^T y) : \|y\|_\infty \leq \lambda\right\}
\]
其中 $f^*(y) = \frac{1}{2}\|y + b\|^2 - \frac{1}{2}\|b\|^2$。

等价于：
\[
\min_y \, \frac{1}{2}\|A^T y - b\|^2 \quad \text{s.t. } \|y\|_\infty \leq \lambda
\]
\end{example}

% 第6章第8节：总结
\section{本章小结}\label{sec:ch6-summary}

\subsection{方法体系总览}

\begin{table}[H]
\centering
\caption{有约束最优化理论方法总结}
\begin{tabular}{@{}>{\raggedright\arraybackslash}p{2.8cm}>{\raggedright\arraybackslash}p{3.2cm}>{\raggedright\arraybackslash}p{3.5cm}>{\raggedright\arraybackslash}p{3.5cm}@{}}
\toprule
\textbf{类别} & \textbf{内容} & \textbf{关键方法} & \textbf{适用场景} \\
\midrule
\textbf{最优性条件}  & KKT一阶条件 & 平稳性+原始/对偶可行性+互补松弛 & 一般NLP \\
 & 二阶条件 & 临界锥+Hessian判定 & 严格局部最优 \\
\midrule
\textbf{约束规范} 
 & LICQ, MFCQ & 积极约束梯度线性无关 & KKT必要性保证 \\
 & Slater条件 & 凸问题严格可行性 & 凸优化强对偶 \\
\midrule
\textbf{对偶理论} 
 & 拉格朗日对偶 & 弱对偶+强对偶定理 & 问题转换、下界估计 \\
 & LP对偶 & 原始-对偶关系 & 线性规划 \\
 & Fenchel对偶 & 共轭函数 & 复合凸优化 \\
\midrule
\textbf{惩罚函数法} 
 & 二次惩罚 & 大参数趋近精确解 & 简单约束 \\
 & 精确惩罚 & 有限参数精确解 & 一般约束 \\
 & 增广拉格朗日 & 乘子更新 & 等式约束 \\
\midrule
\textbf{障碍/内点法} 
 & 对数障碍 & 中心路径 & 不等式约束 \\
 & 原始-对偶 & 牛顿方向+步长控制 & LP, QP, SDP \\
\bottomrule
\end{tabular}
\end{table}

\subsection{KKT条件速查}

\begin{kktbox}[标准形式NLP的KKT条件]
问题：$\min f(x)$ s.t. $g_i(x) \leq 0$，$h_j(x) = 0$

\textbf{平稳性}：$\nabla f(x) + \sum_i \lambda_i \nabla g_i(x) + \sum_j \mu_j \nabla h_j(x) = 0$

\textbf{原始可行}：$g_i(x) \leq 0$，$h_j(x) = 0$

\textbf{对偶可行}：$\lambda_i \geq 0$

\textbf{互补松弛}：$\lambda_i g_i(x) = 0$
\end{kktbox}

\subsection{数值方法选择指南}

\begin{conceptbox}[约束优化方法选择]
\begin{enumerate}
    \item \textbf{等式约束}：
    \begin{itemize}
        \item 少量约束：拉格朗日乘子法 + 牛顿法
        \item 大量约束：增广拉格朗日法
    \end{itemize}
    
    \item \textbf{不等式约束}：
    \begin{itemize}
        \item 凸问题：内点法（推荐）
        \item 非凸问题：SQP（序列二次规划）或积极集法
    \end{itemize}
    
    \item \textbf{线性和二次规划}：
    \begin{itemize}
        \item LP：单纯形法或内点法
        \item QP：积极集法或内点法
    \end{itemize}
    
    \item \textbf{大规模问题}：
    \begin{itemize}
        \item 一阶方法：ADMM, 投影梯度法
        \item 二阶方法：截断牛顿 + 内点
    \end{itemize}
\end{enumerate}
\end{conceptbox}

\subsection{关键定理回顾}

\begin{enumerate}
    \item \textbf{Farkas引理}：两个互斥系统的择一性，KKT理论的基石
    
    \item \textbf{KKT必要条件}：在CQ下，局部最优解必满足KKT条件
    
    \item \textbf{KKT充分条件}：凸优化中，KKT条件等价于全局最优
    
    \item \textbf{弱对偶}：$d^* \leq p^*$ 恒成立
    
    \item \textbf{强对偶}：凸+Slater $\Rightarrow$ $d^* = p^*$
    
    \item \textbf{鞍点定理}：鞍点 $\Leftrightarrow$ 无对偶间隙的最优解对
\end{enumerate}

\subsection{收敛性比较}

\begin{table}[H]
\centering
\caption{数值方法收敛特性}
\begin{tabular}{@{}lccc@{}}
\toprule
\textbf{方法} & \textbf{收敛条件} & \textbf{收敛速度} & \textbf{参数策略} \\
\midrule
二次惩罚 & $\rho \to \infty$ & 线性（慢） & 递增 \\
精确惩罚 & $\rho \geq \bar{\rho}$（有限）& 有限步精确 & 需估计 \\
增广拉格朗日 & $\rho$ 固定或有限 & 超线性 & 乘子更新 \\
对数障碍 & $\mu \to 0$ & 线性（慢）& 递减 \\
原始-对偶内点 & 综合策略 & 超线性 & 自适应 \\
\bottomrule
\end{tabular}
\end{table}

\begin{notebox}[学习建议]
\begin{enumerate}
    \item 首先掌握KKT条件的推导和验证
    \item 理解约束规范的必要性
    \item 掌握对偶理论的弱对偶和强对偶
    \item 了解惩罚函数法和障碍函数法的基本思想
    \item 熟悉内点法的中心路径概念
    \item 能够求解简单的对偶问题
\end{enumerate}
\end{notebox}

\newpage

% ==================== 第七章 约束问题的最优化方法 ====================
